\begin{titlepage}
\centering
{\Large Notas de entrevista\par}
\vspace{1.2cm}
{\huge\bfseries \notetitle\par}
\vspace{0.8cm}
{\large \noteauthors\par}
\vspace{0.3cm}
{\large \notedate\par}
\vspace{1.2cm}
\ifdefempty{\videocoverpath}{}{\includegraphics[width=0.82\textwidth,height=0.45\textheight,keepaspectratio]{\videocoverpath}\par}
\vfill
\begin{tcolorbox}[width=0.9\textwidth, colback=black!2!white, colframe=black!60, sharp corners]
\textbf{Autor o canal del video}: \videochannel\par
\textbf{Fecha de publicación}: \videopublishdate\par
\textbf{Duración del video}: \videoduration\par
\ifdefempty{\videourl}{}{\textbf{Enlace del video}: \href{\videourl}{\nolinkurl{\videourl}}\par}
\textbf{Página del proyecto}: \href{\repourl}{\nolinkurl{\repourl}}\par
\end{tcolorbox}
\vspace{0.3cm}
{\small Esta discusión gira en torno a Agentic training, GUI/mobile agent, memory/tool use, benchmark y la implementación en la industria; el objeto de la discusión no es un modelo único, sino cómo coevolucionan la «capacidad del modelo» y la «ingeniería de sistemas».\par}
\end{titlepage}
\tableofcontents
\newpage

